\section*{अध्याय \ref{ch:b2:mammal-reproduction} --- स्तनधारियों का लैंगिक जनन}

\begin{solution}{exo:b2:mammal-reproduction:1}
शुक्राणुजन यौवनारंभ से जीवन भर विभाजित होते हैं; अंडाणुजन केवल
भ्रूण में। एक \omterm{def:b2:meiosis-heredity:meiosis}{अर्धसूत्री विभाजन} चार शुक्राणु देता है पर एक ही अंडाणु (और
तीन ध्रुवीय पिंड)। किसी शुक्राणु को लगभग 64 दिन और दो सप्ताह का परिपक्वन
लगता है; जबकि कोई अंडाणु जन्म से पहले से लेकर अपने \omterm{def:b2:mammal-reproduction:ovary}{अंडोत्सर्ग} के महीने तक
--- यानी दशकों --- \omterm{prop:b2:meiosis-heredity:stages}{पूर्वावस्था I} में ठहरा रहता है, और फिर निषेचन तक
मध्यावस्था II में। कोई पुरुष प्रतिदिन $10^{8}$ शुक्राणु बनाता है, यानी
जीवन भर में कोई $10^{12}$; जबकि कोई स्त्री लगभग 400 अंडाणु अंडोत्सर्जित करती है।
\end{solution}

\begin{solution}{exo:b2:mammal-reproduction:2}
\omterm{prop:b2:mammal-reproduction:fertilisation}{पारदर्शी पटल} से बँधना; \omterm{prop:b2:mammal-reproduction:fertilisation}{एक्रोसोम अभिक्रिया} और एक रास्ते का
पचाया जाना; झिल्लियों का संलयन और शुक्राणु-केंद्रक का प्रवेश;
कैल्सियम तरंग और \omterm{prop:b2:mammal-reproduction:fertilisation}{वल्कुटीय अभिक्रिया} (पटल का कड़ा होना, उसके
ग्राहियों का लोप: यानी बहुशुक्राणुता का अवरोध); \omterm{def:b2:meiosis-heredity:meiosis}{अर्धसूत्री विभाजन} II का
पूरा होना और दूसरे ध्रुवीय पिंड का निकलना; दोनों पूर्वकेंद्रकों का बनना,
DNA प्रतिकृतिकरण, और पहला समसूत्री विभाजन।
\end{solution}

\begin{solution}{exo:b2:mammal-reproduction:3}
\omterm{def:b2:mammal-reproduction:implantation}{पोषकोरक} (बाहरी परत) \omterm{def:b2:mammal-reproduction:placenta}{अपरा} और झिल्लियाँ बनाता है; और
\omterm{def:b2:mammal-reproduction:implantation}{अंतःकोशिका पिंड} भ्रूण बनाता है। \omterm{def:b2:mammal-reproduction:implantation}{रोपण} निषेचन के बाद दिन 6 से 7
पर आरंभ होता है, ऐसे गर्भाशयी अस्तर में जिसे \omterm{def:b2:mammal-reproduction:ovary}{पीत पिंड} के
प्रोजेस्टेरोन ने तैयार किया हो।
\end{solution}

\begin{solution}{exo:b2:mammal-reproduction:4}
ऑक्सीजन: विसरण। ग्लूकोज़: सुसाध्य अभिगमन। मातृ लाल
कोशिकाएँ: पार नहीं करतीं (दोनों रक्त कभी नहीं मिलते)। IgG: ग्राही-मध्यित
अभिगमन। यूरिया: विसरण, भ्रूण से माता की ओर। ऐल्कोहॉल: स्वतंत्र रूप से
विसरित। जीवाणु: अधिकतर पार नहीं करते (कुछ, जैसे उपदंश और
लिस्टीरिया, करते हैं)।
\end{solution}

\begin{solution}{exo:b2:mammal-reproduction:5}
$m = 3, 1, 0.4, 0.1$: $P = 0.95, 0.63, 0.33, 0.095$। एक-आधे से नीचे तब,
जब $m < \ln 2 = 0.69$ हो, यानी $7\times 10^{7}$ शुक्राणुओं से नीचे। वह
वक्र संतृप्त हो जाता है: लगभग $10^{8}$ से ऊपर अतिरिक्त शुक्राणु लगभग कुछ
नहीं जोड़ते, और उससे नीचे उर्वरता संख्या के लगभग समानुपात में गिरती है।
\end{solution}

\begin{solution}{exo:b2:mammal-reproduction:6}
यौवनारंभ से पहले: 13 वर्षों (\num{4750} दिनों) में $7\times 10^{5}$ खोए:
यानी प्रतिदिन लगभग 150। बाद में: 37 वर्षों (\num{13500}
दिनों) में \num{299000} खोए: प्रतिदिन 22; और \omterm{def:b2:mammal-reproduction:ovary}{अंडोत्सर्ग} $400/\num{13500} = 0.03$
प्रतिदिन। पुटकनाश \omterm{def:b2:mammal-reproduction:ovary}{अंडोत्सर्ग} से सात सौ गुना अधिक है; वह भंडार
\omterm{def:b2:mammal-reproduction:ovary}{अंडोत्सर्ग} से नहीं, मरने से चुकता है।
\end{solution}

\begin{solution}{exo:b2:mammal-reproduction:7}
वयस्क: $(30/27)^{2.7} = 1.33$, $S = 1.33/2.33 = 0.57$। भ्रूणीय:
$(30/19)^{2.7} = 3.4$, $S = 3.4/4.4 = 0.77$। \omterm{def:b2:mammal-reproduction:placenta}{अपरा} के कम दाबों
पर भ्रूणीय रक्त वयस्क रक्त से पाँचवाँ भाग अधिक ऑक्सीजन
भरता है, और उसे ऊतकों में उन दाबों पर उतार देता है जिन पर वयस्क
हीमोग्लोबिन पहले ही आधा ख़ाली हो चुका होता।
\end{solution}

\begin{solution}{exo:b2:mammal-reproduction:8}
$2^{k} = 25$: $k = 4.6$ दुगुनेपन, \qty{9.3}{d}: यानी निषेचन के बाद दिन
16 से 17; और अंतिम रजःस्राव के बाद दिन 30 से 31 --- यानी लगभग उसी दिन
जब रजःस्राव छूटता है।
\end{solution}

\begin{solution}{exo:b2:mammal-reproduction:9}
$3\times 10^{-8}\times 6\times 10^{4}\times 20/10^{-3} =
\qty{36}{mL/min}$, बनाम \qty{25}{mL/min} की माँग: गुंजाइश
आठ से घटकर डेढ़ रह जाती है; वह भ्रूण किसी भी अतिरिक्त भार के नीचे
अल्पऑक्सीजनी हो जाता है, धीरे बढ़ता है और उसे जल्दी जन्माना पड़ सकता है।
\end{solution}

\begin{solution}{exo:b2:mammal-reproduction:10}
(क) \omterm{def:b2:mammal-reproduction:testis}{वृषण} बनते हैं और दोनों हॉर्मोन स्रावित करते हैं; AMH म्यूलरी
नलिकाएँ हटा देता है, पर टेस्टोस्टेरोन काम नहीं कर पाता: वोल्फ़ीय
नलिकाएँ क्षीण हो जाती हैं और बाह्य जननांग मादा होते हैं --- यानी
\omterm{def:b2:mammal-reproduction:testis}{वृषण} वाली और गर्भाशय-रहित एक स्त्री। (ख) टेस्टोस्टेरोन काम करता है,
AMH नहीं: यानी अधिवृषण, शुक्रवाहिका और नर जननांगों वाला कोई नर, और साथ
में गर्भाशय तथा अंडवाहिनियाँ। (ग)
अंडाशय और म्यूलरी व्युत्पन्न (न \omterm{def:b2:mammal-reproduction:testis}{वृषण}, न AMH); और अधिवृक्क
ऐंड्रोजन बाह्य जननांगों को भिन्न-भिन्न मात्रा में नरीकृत कर देते हैं। (घ)
\emph{SRY} \omterm{def:b2:mammal-reproduction:testis}{वृषण} बनाता है, जो नर राह थोप देते हैं: यानी दो
X गुणसूत्रों वाला कोई पुरुष, जो बंध्य है क्योंकि \omterm{def:b2:mammal-reproduction:testis}{शुक्राणुजनन} के Y जीन
अनुपस्थित हैं।
\end{solution}

\begin{solution}{exo:b2:mammal-reproduction:11}
सबसे आम (\qty{68}{\%}) दिन 4 से 8 का बँटवारा है: यानी
\omterm{def:b2:mammal-reproduction:implantation}{पोषकोरक} बन जाने के बाद (इसलिए एक \omterm{def:b2:mammal-reproduction:placenta}{अपरा}) पर उल्ब बनने से पहले
(इसलिए दो थैलियाँ)। दिन 3 से पहले हर आधा अपना \omterm{def:b2:mammal-reproduction:implantation}{पोषकोरक} बनाता है,
इसलिए दो \omterm{def:b2:mammal-reproduction:placenta}{अपराएँ} (\qty{30}{\%}); और दिन 8 के बाद उल्ब
पहले ही बिछ चुका होता है और साझा होता है (\qty{2}{\%})। \omterm{def:b2:mammal-reproduction:implantation}{पोषकोरक} के
प्रतिबद्ध हो जाने के बाद \omterm{def:b2:mammal-reproduction:implantation}{अंतःकोशिका पिंड} का बँटना ही स्पष्टतः
सबसे आसान घटना है।
\end{solution}

\begin{solution}{exo:b2:mammal-reproduction:12}
कोई मार्सूपियल जन्म से पहले थोड़ा ही लगाता है: छोटी गर्भावधि, नन्हा
बच्चा, और ऐसा स्तनपायन जिसे सूखा आने पर थैली का बच्चा गिराकर बीच में ही
रोका जा सकता है, इसलिए प्रति प्रयास माता का जोखिम कम है और
वह तुरंत फिर कोशिश कर सकती है। जबकि कोई अपरायुक्त लंबी
गर्भावधि में भारी निवेश करता है और ऐसा बड़ा, भली-भाँति विकसित बच्चा देता
है जो जन्म पर बेहतर बचता है, पर इसका दाम यह है कि माता महीनों बँधी और
इस बीच असुरक्षित रहती है। उत्पादक, पूर्वानुमेय पर्यावरणों में
अपरायुक्त बच्चों की ऊँची उत्तरजीविता जीतती है; ऑस्ट्रेलिया की अनियमित
जलवायु मार्सूपियल के छोड़ देने के विकल्प को पुरस्कृत करती है, और उसके
अलगाव ने अपरायुक्तों को हाल तक बाहर रखा।
\end{solution}

\begin{solution}{pb:b2:mammal-reproduction:1}
\textbf{1.} $m = 3\times 10^{8}\times 10^{-8} = 3$; $P = 1 - e^{-3} =
0.95$।
\textbf{2.} $1 - e^{-3}(1 + 3) = 0.80$: उस अवरोध के बिना पाँच में चार
निषेचन बहुशुक्राणुक होते और भ्रूण त्रिगुणित।
\textbf{3.} $0.15/5\times 10^{-5} = \qty{3000}{s}$, यानी पचास मिनट;
उन्हें गर्भाशय और अंडवाहिनी के संकुचन ले जाते हैं।
\textbf{4.} $300/3\times 10^{8} = 10^{-6}$: वे अम्लीय योनि में, ग्रीवा
के श्लेष्म में, गर्भाशय के भक्षकाणुओं के हाथों, और ग़लत अंडवाहिनी में
खो जाते हैं।
\textbf{5.} \omterm{def:b2:mammal-reproduction:ovary}{अंडोत्सर्ग} से तीन दिन पहले से एक दिन बाद तक: यानी लगभग
चार दिन।
\textbf{6.} $m = 0.5$, $P = 0.39$; प्रत्याशित चक्र $1/P = 2.5$।
\textbf{7.} शुक्राणु की थोड़ी-सी सूत्रकणिकाएँ, जो मध्यखंड में हैं,
प्रवेश के बाद चिह्नित और नष्ट कर दी जाती हैं; जबकि अंडाणु एक लाख ढोता है।
\textbf{8.} $120/16 = 7.5$: यानी सात विभाजन, $2^{7} = 128$ कोशिकाएँ ---
यानी लगभग सौ।
\textbf{9.} $\log_2 25 = 4.6$ दुगुनेपन, \qty{9.3}{d}: दिन 16 से 17;
और रजःस्राव के बाद दिन 30 से 31।
\textbf{10.} \omterm{def:b2:mammal-reproduction:ovary}{पीत पिंड} \omterm{def:b2:mammal-reproduction:ovary}{अंडोत्सर्ग} के लगभग बारह दिन बाद
क्षीण हो जाता है, जब तक hCG उसे न बचा ले; उसके बिना प्रोजेस्टेरोन गिरता
है, अस्तर झड़ जाता है और वह भ्रूण रजःस्राव के साथ खो जाता है।
\textbf{11.} $\log_2 10^{5} = 16.6$ दुगुनेपन; और सप्ताह 10 तक
\omterm{def:b2:mammal-reproduction:placenta}{अपरा} अपना प्रोजेस्टेरोन स्वयं बनाती है तथा \omterm{def:b2:mammal-reproduction:ovary}{पीत पिंड} की अब
ज़रूरत नहीं रहती।
\textbf{12.} अंडाणु दशकों तक \omterm{prop:b2:meiosis-heredity:stages}{पूर्वावस्था I} में बैठे रहते हैं और
द्विसंयोजियों को थामे कोहेसिन क्षीण हो जाते हैं, इसलिए \omterm{def:b2:meiosis-heredity:meiosis}{अर्धसूत्री विभाजन} I
पर अपृथक्करण मातृ आयु के साथ तेज़ी से चढ़ता है; जबकि शुक्राणु दो महीनों
में ताज़े बनते हैं।
\textbf{13.} \omterm{def:b2:mammal-reproduction:implantation}{पोषकोरक} के बाद (दिन 5) और उल्ब से पहले (दिन
8): यानी एक \omterm{def:b2:mammal-reproduction:placenta}{अपरा}, दो थैलियाँ।
\textbf{14.} $3\times 10^{-8}\times 1.2\times 10^{5}\times 20/3.5\times
10^{-4} = \qty{206}{mL/min}$।
\textbf{15.} $7\times 3.5 = \qty{24.5}{mL/min}$: यानी आठ की गुंजाइश।
\textbf{16.} $500\times 0.16 = \qty{80}{mL/min}$; निष्कर्षण
$24.5/80 = \qty{31}{\%}$।
\textbf{17.} $5\times 3.5\times 1440 = \qty{25}{g/d}$; जो \qty{28}{L}
मातृ रक्त में है; और दैनिक प्रवाह \qty{720}{L}: यानी वह भ्रूण गुज़रते
ग्लूकोज़ का लगभग \qty{4}{\%} लेता है।
\textbf{18.} भ्रूणीय हीमोग्लोबिन \qty{30}{mmHg} पर \qty{77}{\%}
संतृप्त है, उस भ्रूण के पास अधिक लाल कोशिकाएँ और ऊँचा हृद्-निर्गत है,
और उसके ऊतक कम दाब पर काम करने को अनुकूलित हैं --- वह, वस्तुतः,
एवरेस्ट की चोटी पर जीता है।
\textbf{19.} बिना काम करते फेफड़े के गैस-विनिमय, बिना काम करती आँत या
वृक्क के पोषण और उत्सर्जन, और वे हॉर्मोन जो
गर्भावस्था बनाए रखते हैं; पर वह हर विष बाहर नहीं रख सकती --- ऐल्कोहॉल,
निकोटीन और कुछ विषाणु पार कर जाते हैं।
\textbf{20.} ग्रीवा का खिंचाव $\to$ अधश्चेतक $\to$ ऑक्सीटोसिन $\to$
संकुचन $\to$ और खिंचाव, जब तक प्रसव वह उद्दीपन हटा न दे।
पूर्णावधि से पहले प्रोजेस्टेरोन-प्रधान गर्भाशय में ऑक्सीटोसिन
ग्राही थोड़े हैं और अंतराल संधियाँ नहीं, इसलिए संकुचन फैलते नहीं और वह
चक्र बंद नहीं हो पाता।
\textbf{21.} $750\times 2.9 = \qty{2.2}{MJ}$; लागत $2.2/0.8 =
\qty{2.7}{MJ}$।
\textbf{22.} वृद्धि $30\times 6 = \qty{0.18}{MJ}$, अनुरक्षण
$300\times 4 = \qty{1.2}{MJ}$: यानी \qty{1.4}{MJ}, दूध की ऊर्जा का
लगभग दो तिहाई; और बाक़ी क्रियाशीलता, गर्मी और हानियों में जाता है।
\textbf{23.} $270\times 2.7 = \qty{730}{MJ}$, यानी गर्भावस्था का
ढाई गुना।
\textbf{24.} चूसना प्रोलैक्टिन चढ़ाता है और उस हॉर्मोन का स्पंदी
मोचन दबा देता है जो \omterm{def:b2:mammal-reproduction:ovary}{अंडोत्सर्ग} चलाता है, इसलिए \omterm{def:b2:mammal-reproduction:ovary}{अंडोत्सर्ग}
महीनों बाद ही लौटता है; और गर्भनिरोध के बिना यह जन्मों में दो
से तीन वर्ष का अंतर डाल देता है।
\textbf{25.} $P = 0.95$; परीक्षण निषेचन के बाद दिन 17 पर धनात्मक
(चक्र का दिन 31); ऑक्सीजन-गुंजाइश आठ; और एक दिन का दूध माता को
\qty{2.7}{MJ} पड़ता है।
\end{solution}
